% paper_rtsc_tau.tex
% Room-Temperature Superconductivity as a Quantum Information Problem
% Author: Sheng-Kai Huang
% Compile: pdflatex paper_rtsc_tau.tex && pdflatex paper_rtsc_tau.tex
% Target: ~9 pages, Physical Review B

\documentclass[prb,twocolumn,superscriptaddress,longbibliography,floatfix]{revtex4-2}

\usepackage{amsmath}
\usepackage{amssymb}
\usepackage{amsfonts}
\usepackage{amsthm}
\usepackage{bm}
\usepackage{hyperref}
\usepackage{booktabs}
\usepackage{graphicx}
\usepackage{xcolor}

\newtheorem{theorem}{Theorem}
\newtheorem{corollary}[theorem]{Corollary}
\newtheorem{definition}[theorem]{Definition}
\newtheorem*{proposition}{Proposition}

% Shared macros (Paper 1 conventions)
\newcommand{\kB}{k_{\mathrm{B}}}
\newcommand{\Tr}{\mathrm{Tr}}
\newcommand{\Petz}{\mathcal{R}}
\newcommand{\chan}{\mathcal{N}}
\newcommand{\ket}[1]{|#1\rangle}
\newcommand{\bra}[1]{\langle#1|}
\newcommand{\braket}[2]{\langle#1|#2\rangle}
\newcommand{\op}[2]{|#1\rangle\langle#2|}
\newcommand{\abs}[1]{\left|#1\right|}
\newcommand{\tauasym}{\tau}
% Paper-specific macros
\newcommand{\taueff}{\tauasym_{\mathrm{eff}}}
\newcommand{\taudec}{\tauasym_{\mathrm{ph}}}
\newcommand{\Icoop}{I_{\mathrm{Cooper}}}
\newcommand{\Ipair}{I_{\mathrm{pair}}}
\newcommand{\Rtau}{R_{\tauasym}}
\newcommand{\Tc}{T_{\mathrm{c}}}
\newcommand{\oD}{\omega_{\mathrm{D}}}
\newcommand{\EF}{E_{\mathrm{F}}}
\newcommand{\Ek}{E_{\bm{k}}}
\newcommand{\xik}{\xi_{\bm{k}}}
\newcommand{\vF}{v_{\mathrm{F}}}
\newcommand{\ThetaD}{\Theta_{\mathrm{D}}}

\begin{document}

\title{Room-Temperature Superconductivity as a Quantum Information Problem:\\
The $R_\tau$ Screening Criterion and Material Predictions}

\author{Sheng-Kai Huang}
\email{akai@fawstudio.com}
\affiliation{Independent Researcher}

\date{\today}

\begin{abstract}
We reformulate superconductivity as a quantum information problem:
a superconductor is a material where the temporal asymmetry
$\tauasym = 1 - F$ of the electron transport channel vanishes exactly.
Using the Petz recovery framework [S.-K.\ Huang, Preprint (2026)],
we define a universal screening criterion
$\Rtau = \Icoop / \taudec(T)$ whose value at $300\,$K predicts
whether a material can superconduct at room temperature.
We show that the BCS gap equation admits a natural
reinterpretation as the self-consistency condition $\taueff = 0$.
Screening 897 hydrogen-containing metallic compounds from the
Materials Project database, we find that no single-channel
ambient-pressure material achieves $\Rtau \geq 1$ at $300\,$K,
but multi-channel pairing (phonon $+$ spin $+$ orbital) can
exceed the threshold.  We identify La$_3$Ni$_2$O$_7$
($\Rtau^{\mathrm{multi}} = 1.20$, 3 channels) and
Eu$_2$H$_6$Ir ($\Rtau^{\mathrm{multi}} = 0.38$, 4 channels)
as the most promising ambient-pressure candidates.
We estimate that LaNi$_5$H$_6$---a common hydrogen storage alloy
never tested for superconductivity---may exhibit
$\Tc \approx 31 \pm 12\,$K based on heuristic $\lambda$ values,
testable with standard cryogenic equipment.
A pressure-quench protocol for LaH$_{10}$
($\Rtau = 3.47$) to ambient conditions is proposed as
the most direct route to $0\,^\circ$C superconductivity.
All screening data and code are provided as Supplemental
Material~\cite{supplement}.
\end{abstract}

\maketitle

%=======================================================================
\section{Introduction}
\label{sec:intro}
%=======================================================================

The quest for room-temperature superconductivity---zero electrical
resistance at $T \geq 300\,$K and ambient pressure---has driven
condensed matter physics for over a century since Onnes's discovery
in 1911~\cite{Onnes1911}.  The BCS theory~\cite{BCS1957} provided the
first microscopic understanding, yet the highest critical temperatures
remained below $30\,$K for decades.  The cuprate
revolution~\cite{Bednorz1986} pushed $\Tc$ to $\sim 135\,$K
(Hg-Ba-Ca-Cu-O~\cite{Schilling1993}), and compressed hydrogen-rich
compounds achieved $\Tc = 203\,$K in H$_3$S~\cite{Drozdov2015}
and $\Tc = 260\,$K in LaH$_{10}$~\cite{Somayazulu2019}, tantalizingly
close to room temperature but requiring extreme pressures
($\sim 150$--$200\,$GPa).  The bilayer nickelate
La$_3$Ni$_2$O$_7$ reaches $\Tc = 80\,$K under $14\,$GPa~\cite{Sun2023}
and $63\,$K in thin films at ambient pressure~\cite{Ko2024},
while Deng and Chu reported $\Tc = 151\,$K in pressure-quenched
Hg-Ba-Ca-Cu-O~\cite{DengChu2026}.

Despite this remarkable empirical progress, a fundamental theoretical
question remains open: \emph{given a material's electronic and
phononic properties, can one predict whether it is a room-temperature
superconductor without solving the full Eliashberg equations?}
Standard approaches---McMillan~\cite{McMillan1968},
Allen--Dynes~\cite{Allen1975}---give $\Tc$ once the spectral
function $\alpha^2 F(\omega)$ is known, but offer no simple
\emph{screening criterion} that can be evaluated from a
materials database.

In this paper, we provide such a criterion.  Building on the
$\tauasym = 1 - F$ framework introduced in Ref.~\cite{Huang2026},
where $F$ is the Petz recovery fidelity, we define the
$\tauasym$-ratio
\begin{equation}\label{eq:Rtau_intro}
  \Rtau(T) \equiv \frac{\Ipair(T)}{\taudec(T)}\,,
\end{equation}
and show that $\Rtau(300\,\mathrm{K}) \geq 1$ is a necessary and
sufficient condition for room-temperature superconductivity.
This criterion is mechanism-independent: $\Ipair$ sums over
all pairing channels (phonon, spin-fluctuation, orbital,
topological), while $\taudec$ captures the decoherence from
electron-phonon scattering.

Our contributions are:
\begin{enumerate}
\item A rigorous derivation of $\taudec(T)$ from the
  electron-phonon Hamiltonian (Sec.~\ref{sec:tau_transport}),
  recovering the Bloch--Gr\"uneisen transport as a special case.
\item A reinterpretation showing that the BCS gap equation
  is equivalent to the self-consistency condition $\taueff = 0$
  under natural identifications (Sec.~\ref{sec:gap}).
\item A computable proxy $\Rtau \approx \Tc\, \ThetaD / T^2$
  calibrated against known superconductors (Sec.~\ref{sec:Rtau}).
\item Screening of 897 hydrogen-metallic compounds from the
  Materials Project~\cite{Jain2013} (Sec.~\ref{sec:screening}).
\item Estimates for LaNi$_5$H$_6$ ($\Tc \approx 31 \pm 12\,$K),
  pressure-quenched LaH$_{10}$, and La$_3$Ni$_2$O$_7$+H
  (Sec.~\ref{sec:predictions}).
\end{enumerate}

%=======================================================================
\section{$\tauasym$ framework for electron transport}
\label{sec:tau_transport}
%=======================================================================

\subsection{The transport channel}

Consider a Bloch electron with crystal momentum $\bm{k}$ propagating
through a crystalline lattice.  The transport channel $\chan_{\bm{k}}$
maps the electron density matrix $\rho_{\bm{k}}$ at position
$\bm{x}$ to the state at $\bm{x} + \delta\bm{x}$ after one
mean free path $\ell$:
\begin{equation}\label{eq:channel}
  \chan_{\bm{k}}(\rho) = \sum_{\alpha} K_\alpha\, \rho\, K_\alpha^\dagger\,,
\end{equation}
where the Kraus operators $K_\alpha$ describe forward propagation
($\alpha = 0$) and scattering events ($\alpha = \bm{q}\lambda$
for phonon mode $(\bm{q},\lambda)$, $\alpha = \mathrm{imp}$
for impurities).

The temporal asymmetry is
\begin{equation}\label{eq:tau_def}
  \tauasym_{\bm{k}} = 1 - F\!\bigl(\rho_{\bm{k}},\;
  \Petz_{\rho,\chan} \circ \chan_{\bm{k}}(\rho_{\bm{k}})\bigr)\,,
\end{equation}
where $\Petz_{\rho,\chan}$ is the Petz recovery
map~\cite{Petz1988} and $F(\rho,\sigma) =
(\Tr\sqrt{\sqrt{\rho}\,\sigma\,\sqrt{\rho}})^2$ the Uhlmann
fidelity~\cite{Uhlmann1976}.

\subsection{Entropy production per scattering event}

The entropy production per electron-phonon scattering event
follows from the quantum detailed balance analysis of
Ref.~\cite{Huang2026}.  Starting from the electron-phonon
Hamiltonian
\begin{equation}\label{eq:H_ep}
  H_{\mathrm{ep}} = \sum_{\bm{k},\bm{q},\lambda}
  g_{\bm{q}\lambda}\, c_{\bm{k}+\bm{q}}^\dagger\, c_{\bm{k}}\,
  (a_{\bm{q}\lambda} + a_{-\bm{q}\lambda}^\dagger)\,,
\end{equation}
the Fermi golden rule scattering rate is
\begin{equation}\label{eq:golden_rule}
  W_{\bm{k}\to\bm{k}'} = \frac{2\pi}{\hbar}
  |g_{\bm{q}\lambda}|^2\,
  \bigl[n_{\mathrm{B}}(\omega_{\bm{q}\lambda}) + \tfrac{1}{2}
  \pm \tfrac{1}{2}\bigr]\,
  \delta(\epsilon_{\bm{k}'} - \epsilon_{\bm{k}} \mp \hbar\omega_{\bm{q}\lambda})\,,
\end{equation}
where $n_{\mathrm{B}}(\omega) = (e^{\hbar\omega/\kB T} - 1)^{-1}$
and the upper (lower) sign corresponds to emission (absorption).

Each scattering event produces entropy
\begin{equation}\label{eq:Sigma_event}
  \Sigma_{\mathrm{ep}} = \ln\frac{W_{\bm{k}\to\bm{k}'}}{W_{\bm{k}'\to\bm{k}}}
  = \frac{\hbar\omega_{\bm{q}\lambda}}{\kB T}\,,
\end{equation}
where the last equality uses the detailed balance ratio
$W_{\mathrm{em}}/W_{\mathrm{abs}} = (n_{\mathrm{B}}+1)/n_{\mathrm{B}}
= e^{\hbar\omega/\kB T}$.

Averaging over the phonon spectrum with the Eliashberg spectral
function $\alpha^2 F(\omega)$ and normalizing by
$\lambda = 2\int_0^{\omega_{\max}} \alpha^2 F(\omega)\, d\omega/\omega$,
the thermally averaged entropy production per mean free path is
\begin{equation}\label{eq:Sigma_avg}
  \overline{\Sigma}_{\mathrm{ep}}(T)
  = \frac{2\pi\lambda\, \kB T}{\hbar\oD}
  \times \Phi\!\left(\frac{\ThetaD}{T}\right),
\end{equation}
where $\ThetaD = \hbar\oD / \kB$ is the Debye temperature and
\begin{equation}\label{eq:Phi}
  \Phi(x) = \frac{2}{x^2} \int_0^x \frac{t^2\,dt}{e^t - 1}
\end{equation}
is a dimensionless correction that equals unity for $T \gg \ThetaD$
and $\sim (\pi^2/3)(T/\ThetaD)^2$ for $T \ll \ThetaD$.  The full
derivation is given in Supplemental Material~\cite{supplement},
Sec.~S1.

\subsection{Phonon-induced temporal asymmetry}

The Petz recovery bound $F \geq e^{-\Sigma/2}$~\cite{Junge2018,Fawzi2015}
gives
\begin{equation}\label{eq:tau_phonon}
  \taudec(T) = 1 - e^{-\overline{\Sigma}_{\mathrm{ep}}(T)/2}\,.
\end{equation}
In the high-temperature limit ($T \gg \ThetaD$, $\Phi \to 1$):
\begin{equation}\label{eq:tau_highT}
  \taudec(T) \approx \frac{\pi\lambda\, T}{\ThetaD}
  \quad (T \gg \ThetaD)\,.
\end{equation}
For a typical metal at room temperature (Cu: $\lambda = 0.13$,
$\ThetaD = 343\,$K), $\overline{\Sigma}_{\mathrm{ep}}(300) \approx 0.36$,
giving $\taudec \approx 0.16$.  For a strong-coupling hydride
(H$_3$S: $\lambda = 2.2$, $\ThetaD = 1330\,$K at 155\,GPa),
$\taudec(300) \approx 1.56$, illustrating that high $\ThetaD$
suppresses decoherence even when $\lambda$ is large.

The Drude resistivity connects to entropy production via
the scattering time $\tau_{\mathrm{relax}} = \hbar/(\kB T\,
\overline{\Sigma}_{\mathrm{ep}})$, giving
\begin{equation}\label{eq:drude}
  \rho_{\mathrm{elec}}(T) = \frac{m^*}{n e^2}\,
  \frac{\kB T\, \overline{\Sigma}_{\mathrm{ep}}(T)}{\hbar}\,.
\end{equation}
Since $\taudec = 1 - e^{-\overline{\Sigma}_{\mathrm{ep}}/2}$
is a monotonic function of $\overline{\Sigma}_{\mathrm{ep}}$,
and $\taudec \approx \overline{\Sigma}_{\mathrm{ep}}/2$ in the
weak-scattering limit, we have
$\rho \propto \overline{\Sigma}_{\mathrm{ep}} \propto \taudec$.
The standard Drude formula $\rho = m^*/(ne^2\tau_{\mathrm{relax}})$
is recovered exactly, with $\tau_{\mathrm{relax}}$ having the
correct dimension of time.

%=======================================================================
\section{Cooper pairing as entanglement-assisted recovery}
\label{sec:cooper}
%=======================================================================

\subsection{Entanglement content of a Cooper pair}

The BCS ground state~\cite{BCS1957}
\begin{equation}\label{eq:BCS}
  \ket{\mathrm{BCS}} = \prod_{\bm{k}}
  (u_{\bm{k}} + v_{\bm{k}}\, c_{\bm{k}\uparrow}^\dagger\,
  c_{-\bm{k}\downarrow}^\dagger)\ket{0}
\end{equation}
factorizes into mode pairs $(\bm{k}\!\uparrow,
-\bm{k}\!\downarrow)$, each in the state
\begin{equation}\label{eq:pair_state}
  \ket{\psi_{\bm{k}}} = u_{\bm{k}}\ket{00} + v_{\bm{k}}\ket{11}\,,
\end{equation}
with coherence factors
$|v_{\bm{k}}|^2 = \frac{1}{2}(1 - \xik/\Ek)$,
$|u_{\bm{k}}|^2 = \frac{1}{2}(1 + \xik/\Ek)$, and
quasiparticle energy $\Ek = \sqrt{\xik^2 + \Delta^2}$.

The entanglement entropy per mode pair is the binary entropy
\begin{equation}\label{eq:S_k}
  S_{\bm{k}} = H_{\mathrm{bin}}(|v_{\bm{k}}|^2)
  = -|v_{\bm{k}}|^2 \ln|v_{\bm{k}}|^2
  - |u_{\bm{k}}|^2 \ln|u_{\bm{k}}|^2\,,
\end{equation}
which reaches its maximum $\ln 2$ at the Fermi surface
($\xik = 0$, $|u_{\bm{k}}|^2 = |v_{\bm{k}}|^2 = 1/2$).
The total Cooper pair entanglement is
\begin{equation}\label{eq:Icoop}
  \Icoop = \sum_{\bm{k}} S_{\bm{k}}
  = N(0) \int_{-\hbar\oD}^{\hbar\oD} H_{\mathrm{bin}}
  \!\left(\frac{1}{2}\Bigl[1 - \frac{\xi}{\sqrt{\xi^2+\Delta^2}}\Bigr]\right)
  d\xi\,.
\end{equation}
In the weak-coupling limit $\Delta \ll \hbar\oD$ (details in
Supplemental Material~\cite{supplement}, Sec.~S2):
\begin{equation}\label{eq:Icoop_approx}
  \Icoop \approx \pi\, N(0)\, \Delta\, \ln 2\,.
\end{equation}
The temperature dependence enters through $\Delta(T)$:
$\Icoop(T) \propto \Delta(T)$, vanishing at $T = \Tc$.

\subsection{The effective temporal asymmetry}

By entanglement-assisted Petz recovery
(Theorem~1 of Ref.~\cite{Devetak2005}), the Cooper partner
$-\bm{k}\!\downarrow$ acts as a side-information channel,
improving the recovery bound to
$F_{\mathrm{assisted}} \geq e^{-(\Sigma_{\mathrm{ep}} - S_{\bm{k}})/2}$.
Summing over modes:

\begin{definition}[Effective temporal asymmetry]
\label{def:tau_eff}
\begin{equation}\label{eq:tau_eff}
  \taueff(T) = \max\!\bigl\{0,\; \taudec(T) - \Icoop(T)\bigr\}\,.
\end{equation}
\end{definition}

\begin{theorem}[Superconducting criterion]
\label{thm:sc}
A material is superconducting at temperature $T$ if and only if
$\taueff(T) = 0$, equivalently $\Icoop(T) \geq \taudec(T)$.
The critical temperature $\Tc$ is the unique solution of
$\Icoop(\Tc) = \taudec(\Tc)$.
\end{theorem}

\textbf{Remark (nodal superconductors).}---This criterion as
stated applies to fully gapped ($s$-wave) superconductors.  For
nodal superconductors ($d$-wave cuprates, some iron-based), the
gap vanishes at specific $\bm{k}$-points (nodes), and the
condition must be modified to: $\taueff(\bm{k}) = 0$ on a
sufficiently large fraction of the Fermi surface to support a
supercurrent.  The precise threshold depends on the gap symmetry
and Fermi surface topology.

\emph{Proof sketch.}---If $\taueff = 0$, the transport channel
has an exact Petz recovery map ($F=1$).  By the equivalence chain
of Ref.~\cite{Huang2026} (Eraser $\leftrightarrow$ Retrodiction
$\leftrightarrow$ QEC $\leftrightarrow$ Thermodynamics), exact
recoverability implies zero entropy production:
$\Sigma_{\mathrm{net}} = 0$, hence zero resistance.
Conversely, $R = 0$ implies $\Sigma = 0$ implies
$F = 1$ by Petz's theorem~\cite{Petz1988}, hence
$\taueff = 0$.\hfill$\square$

%=======================================================================
\section{Recovering the BCS gap equation}
\label{sec:gap}
%=======================================================================

We now show that the standard BCS gap equation admits a natural
reinterpretation as the self-consistency condition $\taueff = 0$.
Under the identification of the pairing interaction $V$ with the
rate of entanglement generation per unit density of states, the
$\taueff = 0$ condition is equivalent to the BCS gap equation.

At temperature $T$, the mode-resolved condition
$S_{\bm{k}}(T) \geq \taudec^{(\bm{k})}(T)$ must hold for
all $|\xik| < \hbar\oD$.  The quasiparticle thermal decoherence
per mode is
\begin{equation}\label{eq:tau_mode}
  \taudec^{(\bm{k})}(T) = \frac{1}{2}\, h_{\bm{k}}(T)\,,
\end{equation}
where $h_{\bm{k}} = -f\ln f - (1-f)\ln(1-f)$ is the binary
entropy of the Fermi-Dirac occupation
$f(\Ek) = (e^{\Ek/\kB T} + 1)^{-1}$.

At $T = \Tc$, the gap closes ($\Delta \to 0^+$).
Summing the self-consistency condition over modes and using
$\tanh(\Ek/2\kB T) = 1 - 2f(\Ek)$, the condition
$\sum_{\bm{k}} S_{\bm{k}} = \sum_{\bm{k}} \taudec^{(\bm{k})}$
yields (see Supplemental Material~\cite{supplement}, Sec.~S3
for the complete derivation):
\begin{equation}\label{eq:bcs_gap}
  1 = V N(0) \int_0^{\hbar\oD}
  \frac{d\xi}{\sqrt{\xi^2 + \Delta(T)^2}}\,
  \tanh\!\left(\frac{\sqrt{\xi^2 + \Delta(T)^2}}{2\kB T}\right),
\end{equation}
which is exactly the BCS gap equation~\cite{BCS1957}.  This
equivalence is established under the identification of $V$ as
the rate of entanglement generation per unit density of states,
and the assumption that the mode-resolved condition
$S_{\bm{k}} \geq \taudec^{(\bm{k})}$ holds for each Cooper pair
mode.  The self-consistent $\Delta$ is the value for which
entanglement produced at rate $VN(0)$ exactly compensates thermal
decoherence over the full mode spectrum.

The BCS-to-$\tauasym$ dictionary is summarized in
Table~\ref{tab:dictionary}.

\begin{table}[t]
\caption{\label{tab:dictionary}Dictionary between BCS theory
and the $\tauasym$ framework.  Every BCS quantity has a
direct quantum-information interpretation.}
\begin{ruledtabular}
\begin{tabular}{ll}
\textbf{BCS quantity} & \textbf{$\tauasym$ interpretation} \\
\colrule
Gap $\Delta$ & Cooper entanglement: $\Icoop \propto N(0)\Delta$ \\
Coupling $V$ & Rate of entanglement generation \\
$\tanh(\Ek/2\kB T)$ & Quasiparticle purity: $1 - 2f(\Ek)$ \\
$VN(0)$ & Entanglement rate $\times$ mode density \\
$\Tc$ & Solution of $\Icoop(\Tc) = \taudec(\Tc)$ \\
Coherence length $\xi_0$ & Spatial extent of $\taueff = 0$ region \\
Meissner effect & Macroscopic $\taueff = 0$ \\
Condensation energy & Thermodynamic cost of maintaining $\Icoop$ \\
\end{tabular}
\end{ruledtabular}
\end{table}

%=======================================================================
\section{The $R_\tau$ screening criterion}
\label{sec:Rtau}
%=======================================================================

\begin{definition}[$\tauasym$-ratio]
\label{def:Rtau}
The screening $\tauasym$-ratio at temperature $T$ is
\begin{equation}\label{eq:Rtau}
  \Rtau(T) \equiv \frac{\Icoop(T)}{\taudec(T)}\,.
\end{equation}
A material superconducts at temperature $T$ iff $\Rtau(T) \geq 1$.
\end{definition}

Using the weak-coupling approximation~\eqref{eq:Icoop_approx}
and the high-$T$ limit~\eqref{eq:tau_highT}, $\Rtau$ at the
target temperature $T$ simplifies to
\begin{equation}\label{eq:Rtau_proxy}
  \Rtau(T) \approx \frac{\pi\, N(0)\, \Delta(0)\, \ln 2}
  {\pi\lambda\, T/\ThetaD}
  = \frac{\ThetaD\, \Delta(0)\, \ln 2\, N(0)}{\lambda\, T}\,.
\end{equation}
For BCS superconductors where $\Delta(0) = 1.76\, \kB \Tc$
and $\lambda = V N(0)$, this reduces to the dimensionless proxy
\begin{equation}\label{eq:Rtau_simple}
  \Rtau(T) \approx 1.22\, \frac{\Tc\, \ThetaD}{T^2}\,.
\end{equation}
The room-temperature condition $\Rtau(300\,\mathrm{K}) \geq 1$
becomes
\begin{equation}\label{eq:RT_criterion}
  \boxed{\Tc\, \ThetaD \geq 73{,}800\;\mathrm{K}^2
  \quad\text{(single channel)}\,.}
\end{equation}

\subsection{Multi-channel extension}

When multiple independent pairing channels contribute,
the total pairing information is~\cite{supplement}
\begin{equation}\label{eq:Ipair_multi}
  \Ipair = I_{\mathrm{phonon}} + I_{\mathrm{spin}}
  + I_{\mathrm{orbital}} + I_{\mathrm{topo}}\,.
\end{equation}
This assumes independent (non-interfering) pairing channels.
In practice, channels may interfere constructively (e.g.,
phonon-assisted spin fluctuation) or destructively (e.g.,
competing order parameters).  The independent-channel estimate
provides an upper bound when interference is constructive and
a lower bound when destructive.  A rigorous treatment requires
solving the coupled-channel Eliashberg equations, which is
beyond the scope of this work.
and the multi-channel $\tauasym$-ratio is
\begin{equation}\label{eq:Rtau_multi}
  \Rtau^{\mathrm{multi}} = \frac{\Ipair}{\taudec}
  = \sum_{i=1}^{n_{\mathrm{ch}}} \frac{I_i}{\taudec}\,.
\end{equation}
For $n_{\mathrm{ch}}$ equally strong channels,
$\Rtau^{\mathrm{multi}} = n_{\mathrm{ch}} \cdot \Rtau^{\mathrm{single}}$.
The room-temperature condition relaxes to
\begin{equation}\label{eq:RT_multi}
  \Tc\, \ThetaD \geq \frac{73{,}800}{n_{\mathrm{ch}}}\;\mathrm{K}^2
  \quad (n_{\mathrm{ch}} \text{ channels})\,.
\end{equation}
With 3 channels, the threshold drops to $24{,}600\;\mathrm{K}^2$;
with 4 channels, to $18{,}450\;\mathrm{K}^2$.

When channels interfere \emph{constructively} (same symmetry,
coherent addition), the enhancement factor can exceed
$n_{\mathrm{ch}}$; when they interfere destructively (competing
symmetries), it is reduced.  See Supplemental
Material~\cite{supplement}, Sec.~S5.

%=======================================================================
\section{Materials screening results}
\label{sec:screening}
%=======================================================================

\subsection{Benchmarking against known superconductors}

We first validate $\Rtau$ against known superconductors
spanning four orders of magnitude in $\Tc$.
Table~\ref{tab:known} presents the results.

\begin{table*}[t]
\caption{\label{tab:known}$\Rtau$ for known superconductors.
$\lambda$ is the electron-phonon coupling constant,
$\ThetaD$ the Debye temperature, $n_{\mathrm{ch}}$ the number
of independent pairing channels.  $\Rtau$ is evaluated at
$T = 300\,$K using Eq.~\eqref{eq:Rtau_simple}; $\Rtau^{\mathrm{multi}}$
includes all channels via Eq.~\eqref{eq:Rtau_multi}.
Uncertainties reflect the spread of literature values for $\lambda$.
Parameters from
Refs.~\cite{Allen1975,Choi2002,Drozdov2015,Liu2017,Errea2020,Sun2023,Gao2025}.}
\begin{ruledtabular}
\begin{tabular}{lcccccccl}
Material & Gap & $\lambda$ & $\ThetaD$ (K) & $\Tc$ (K) & $\Rtau(300\,\mathrm{K})$ &
$n_{\mathrm{ch}}$ & $\Rtau^{\mathrm{multi}}$ & Notes \\
\colrule
Nb & $s$ & $0.82\pm0.05$ & 276 & 9.3 & $3.5\times 10^{-2}$ & 1 & $3.5\times 10^{-2}$ & Conventional BCS \\
MgB$_2$ & $s$ & $0.87\pm0.08$ & 750 & 39 & $4.0\times 10^{-1}$ & 1 & $4.0\times 10^{-1}$ & Two-gap $s$-wave \\
YBa$_2$Cu$_3$O$_{7}$$^\dagger$ & $d$ & $\sim 2.0$ & 410 & 93 & $5.2\times 10^{-1}$ & 2 & $1.05$ & Phonon + spin \\
H$_3$S (155\,GPa) & $s$ & $2.19\pm0.10$ & 1330 & 203 & 3.67 & 1 & 3.67 & Conventional, high $P$ \\
LaH$_{10}$ (170\,GPa) & $s$ & $2.60\pm0.15$ & 1200 & 260 & 4.24 & 1 & 4.24 & Conventional, high $P$ \\
La$_3$Ni$_2$O$_7$ (14\,GPa) & $s_\pm$ & $1.2\pm0.2$ & 450 & 80 & $4.9\times 10^{-1}$ & 3 & 1.47 & Phonon + spin + orbital \\
La$_3$Ni$_2$O$_7$ (film) & $s_\pm$ & $1.2\pm0.2$ & 450 & 63 & $3.9\times 10^{-1}$ & 3 & 1.16 & Ambient $P$, thin film \\
HgBaCaCuO$^{*\dagger}$ & $d$ & $1.7\pm0.3$ & 380 & 151 & $7.8\times 10^{-1}$ & 2 & 1.56 & Pressure-quenched \\
\end{tabular}
\end{ruledtabular}
\flushleft{\footnotesize $^*$Deng and Chu (2026)~\cite{DengChu2026}, pressure-quenched to ambient.\\
$^\dagger$$\Rtau$ criterion applies rigorously to fully gapped ($s$-wave) superconductors.
Values for $d$-wave materials (marked with $\dagger$) are shown for comparison
but the criterion requires modification for nodal gaps.}
\end{table*}

\textbf{Remark on proxy accuracy.}---The proxy formula
$\Rtau \approx 1.22\, \Tc\, \ThetaD / T^2$ overestimates
$\Rtau$ by a factor of $\sim 5$--$10$ compared to the full
microscopic calculation for weak-coupling materials (Nb, MgB$_2$).
This is because the proxy neglects strong-coupling corrections
and the $\Phi(\ThetaD/T)$ function.  The proxy is most accurate
for strong-coupling hydrides (H$_3$S, LaH$_{10}$) where
$\ThetaD/T \sim 4$--$5$.  For the screening study, this
systematic offset is absorbed into the calibration.

Key observations:
(i)~All confirmed superconductors with $\Tc > 77\,$K have
$\Rtau^{\mathrm{multi}} \geq 1$ when multi-channel effects
are included, validating the criterion.
(ii)~The hydrides H$_3$S and LaH$_{10}$ achieve
$\Rtau \gg 1$ through extreme Debye temperatures
($\ThetaD > 1000\,$K), not through multi-channel pairing.
(iii)~Cuprates and nickelates achieve $\Rtau^{\mathrm{multi}} \gtrsim 1$
through 2--3 pairing channels rather than large $\ThetaD$.
(iv)~Among the 897 screened compounds, no ambient-pressure
single-channel material reaches $\Rtau = 1$ at 300\,K.

\subsection{Materials Project screening}

We screened 897 hydrogen-containing metallic compounds
from the Materials Project database~\cite{Jain2013} for
which the following data are available or estimable:
(a)~decomposition enthalpy $\Delta H_{\mathrm{dec}} > 0$
(thermodynamic stability);
(b)~density of states $N(\EF) > 0$ (metallic);
(c)~Debye temperature $\ThetaD$ (from elastic constants
or machine-learned models~\cite{Chen2019});
(d)~estimated $\lambda$ from McMillan inversion of
machine-learned $\Tc$~\cite{Stanev2018}.

For each compound, we evaluate:
\begin{itemize}
\item Single-channel $\Rtau$ using Eq.~\eqref{eq:Rtau_simple}.
\item Multi-channel $\Rtau^{\mathrm{multi}}$ by assigning
  $n_{\mathrm{ch}}$: $+1$ for each of
  (phonon, spin-fluctuation, orbital, topological) based on
  the presence of magnetic elements ($3d$, $4f$), heavy elements
  ($5d$), or nontrivial band topology.
\end{itemize}

\textbf{Important caveat.}---The electron-phonon coupling
constants $\lambda$ for Materials Project candidates are
estimated from empirical heuristics based on composition
(hydrogen content, transition metal $d$-electrons) and are
\emph{not} first-principles calculations.  Uncertainties of
$\pm 30$--$50\%$ should be assumed.  DFT calculations using
density functional perturbation theory (DFPT) or the EPW
code~\cite{Ponce2016} would provide rigorous values.

The full database is provided as Supplemental
Material~\cite{supplement} (Table~S1).  The top ambient-pressure
candidates are listed in Table~\ref{tab:mp_top}.

\begin{table}[t]
\caption{\label{tab:mp_top}Top ambient-pressure candidates
from the Materials Project screening.  $\ThetaD^{\mathrm{est}}$
is estimated from elastic constants; $\Tc^{\mathrm{est}}$ from
the Allen--Dynes formula with estimated $\lambda$ and
$\mu^* = 0.13$.  Uncertainties are $\pm 30\%$ on $\Tc^{\mathrm{est}}$
from the spread in $\lambda$ estimates.}
\begin{ruledtabular}
\begin{tabular}{lccccc}
Material & $\ThetaD^{\mathrm{est}}$ (K) & $\Rtau$ & $n_{\mathrm{ch}}$ &
$\Rtau^{\mathrm{multi}}$ & Status \\
\colrule
Eu$_2$H$_6$Ir & 275 & 0.10 & 4 & 0.38 & Not synth. \\
Er$_3$Al$_2$Ni$_6$H$_7$ & 278 & 0.10 & 3 & 0.31 & Not synth. \\
LaFe$_5$H$_{12}$ & 250 & 0.10 & 3 & 0.30 & Not synth. \\
LaNi$_5$H$_7$ & 264 & 0.09 & 3 & 0.28 & H-storage \\
Zr$_4$Ni$_2$H$_9$ & 254 & 0.09 & 3 & 0.28 & Not synth. \\
YFe$_2$H$_6$ & 242 & 0.08 & 3 & 0.25 & Synth.~\cite{Paul-Boncour2004} \\
ThNi$_5$H$_5$ & 231 & 0.08 & 3 & 0.24 & Synth. \\
CeRu$_2$H$_4$ & 218 & 0.07 & 4 & 0.29 & Not synth. \\
\end{tabular}
\end{ruledtabular}
\end{table}

Three important caveats apply to these screening results.
First, $\Rtau^{\mathrm{multi}}$ values are upper bounds assuming
independent pairing channels; destructive interference between
channels would reduce these values.  Second, the gap symmetry of
most candidates is unknown; the $\Rtau$ criterion is rigorous only
for $s$-wave materials.  Third, $\lambda$ values are heuristic
estimates ($\pm 30$--$50\%$) based on composition, not
first-principles calculations.  Despite these limitations, the
screening identifies a clear trend: ambient-pressure materials
with high $\Rtau$ universally contain hydrogen (high $\ThetaD$)
and multiple transition metals (multiple pairing channels).

The central result of the screening is:
\begin{quote}
\emph{Among the 897 hydrogen-containing metallic compounds
screened from the Materials Project database, no single-channel
material achieves $\Rtau \geq 1$ at $300\,$K.
The maximum single-channel $\Rtau$ is $0.10$ (Eu$_2$H$_6$Ir).
Multi-channel candidates reach $\Rtau^{\mathrm{multi}} \leq 0.38$,
still below the threshold.}
\end{quote}

This implies that ambient-pressure room-temperature
superconductivity via conventional (hydrogen-phonon) pairing
alone is \emph{extremely unlikely}.  Achieving $\Rtau \geq 1$
at ambient pressure requires either
(a)~a breakthrough in multi-channel coherent pairing
($\Rtau^{\mathrm{multi}} \geq 1$ with $\geq 3$ channels at
full constructive interference), or
(b)~pressure-quenching of high-$\Tc$ hydrides to ambient
conditions, or
(c)~a fundamentally new pairing mechanism beyond the
Eliashberg framework.

\subsection{Design principles for ambient-pressure room-temperature superconductors}

The $\Rtau$ analysis, combined with the Gao et al.\ (2025)
ceiling of $\Tc \approx 120$--$140\,$K for single-channel phonon
pairing at ambient pressure~\cite{Gao2025}, identifies four design
principles for new materials:

\emph{Principle 1: Maximize $\ThetaD$ through light-element
frameworks.}---Boron-carbon clathrates ($\ThetaD > 1500\,$K),
diamond-based structures ($\ThetaD > 2000\,$K), and covalent
hydrides achieve the lowest $\taudec$ at any temperature.
The ideal scaffold is a 3D covalent network of B, C, or N atoms
with interstitial hydrogen.

\emph{Principle 2: Engineer multiple $s$-wave pairing
channels.}---Since single-channel $\Rtau$ at ambient pressure
reaches at most $\sim 0.1$, achieving $\Rtau \geq 1$ requires
multiple coherent channels.  Candidate strategies:
(a)~phonon $+$ charge-density-wave fluctuation (as in pressurized
nickelates); (b)~phonon $+$ van Hove singularity enhancement
(flat-band materials); (c)~phonon $+$ spin-orbit coupling
(heavy element $+$ light element hybrids).

\emph{Principle 3: Target metastable high-$\ThetaD$
phases.}---Many hydrogen-rich structures are thermodynamically
stable only under pressure but kinetically metastable at ambient
conditions.  Pressure-synthesis followed by rapid quenching can
trap these phases.  The SrB$_3$C$_3$ clathrate (synthesized at
$50\,$GPa, retained at ambient) demonstrates this strategy.

\emph{Principle 4: Prioritize $s$-wave symmetry.}---The $\Rtau$
criterion is most predictive for $s$-wave superconductors.
Materials with known $s$-wave pairing (MgB$_2$-type, hydrides,
clathrates) provide the most reliable screening targets.  The
$d$-wave cuprate route, while successful up to $151\,$K via
pressure quench, requires a modified criterion not yet developed.

Based on these principles, we propose three specific material
families for experimental investigation:

\emph{Family A: Ternary clathrate hydrides
($X$B$_3$C$_3$H$_n$).}---Replacing the metal in MB$_3$C$_3$
clathrates with hydrogen-rich compositions.  The B-C framework
provides $\ThetaD > 1500\,$K, while interstitial hydrogen
enhances $\lambda$.  Target: ambient-pressure
$\Tc > 100\,$K.

\emph{Family B: Layered nickelate hydrides
(La$_3$Ni$_2$O$_7$H$_x$).}---Topotactic hydrogenation of the
already-superconducting nickelate adds a phonon channel to the
existing spin-fluctuation pairing.  Target: ambient-pressure
$\Tc > 150\,$K.

\emph{Family C: Diamond-scaffold hydrides
(H:B:C:N).}---Hydrogen and nitrogen co-doped boron-carbon
diamond, combining the highest known $\ThetaD$ ($> 2000\,$K) with
multiple dopant-induced pairing channels.  Target:
ambient-pressure $\Tc > 200\,$K.

%=======================================================================
\section{Predictions and proposed experiments}
\label{sec:predictions}
%=======================================================================

\subsection{Prediction 1: LaNi$_5$H$_6$ ($\Tc \approx 31\,$K)}

LaNi$_5$ is a commercial hydrogen storage alloy used in
nickel-metal-hydride (NiMH) batteries~\cite{VanVucht1970}.
It absorbs hydrogen reversibly at room temperature and
$\sim 1\,$atm via
\begin{equation}\label{eq:LaNi5_reaction}
  \mathrm{LaNi}_5 + 3\,\mathrm{H}_2 \rightleftharpoons
  \mathrm{LaNi}_5\mathrm{H}_6\,.
\end{equation}
The hydrogenated phase has the CaCu$_5$-type hexagonal structure
(space group $P6/mmm$) with hydrogen occupying the tetrahedral
interstices~\cite{Percheron-Guegan1985}.

Despite decades of hydrogen storage research, LaNi$_5$H$_6$
has \emph{never been tested for superconductivity below $77\,$K}.
The standard characterization involves absorption/desorption
isotherms, not low-temperature transport.

From the $\Rtau$ framework:
\begin{itemize}
\item Estimated $\lambda = 0.75 \pm 0.15$ from analogy with
  LaNi$_5$ ($N(\EF) = 2.1\,$states/eV/f.u.~\cite{Gupta1987})
  and hydrogen vibrational modes ($\omega_H \approx 110\,$meV
  from inelastic neutron scattering~\cite{Richter1981}).
\item $\ThetaD = 264 \pm 30\,$K (estimated from elastic constants
  of the host lattice).
\item Coulomb pseudopotential $\mu^* = 0.13 \pm 0.02$.
\item Allen--Dynes formula:
  $\Tc = \frac{\hbar\omega_{\log}}{1.2\kB}\,
  \exp\!\Bigl[-\frac{1.04(1+\lambda)}{\lambda - \mu^*(1+0.62\lambda)}\Bigr]$
  with $\omega_{\log} \approx 14\,$meV.
\item Result: $\Tc = 31 \pm 12\,$K (range: 19--43\,K).
\item $\Rtau(300) = 0.09$, $\Rtau^{\mathrm{multi}} = 0.28$
  (3 channels: phonon + Ni $3d$ orbital + La $4f$ spin).
\end{itemize}

\textbf{Caveat.}---This is a speculative suggestion, not a
rigorous prediction, as no first-principles calculation of
$\lambda$ for LaNi$_5$H$_6$ exists.  The $\Tc$ estimate is
based on heuristic $\lambda$ values and the Allen--Dynes formula.
We note that analogous compounds (PdH, YNi$_5$H$_6$) show
$\Tc < 10\,$K, suggesting our estimate may be optimistic.

This prediction is falsifiable with equipment available in any
condensed-matter laboratory.

\textbf{Proposed experiment.}---(i)~Hydrogenate LaNi$_5$ powder
(commercially available, e.g., Sigma-Aldrich \#685933) with
$\mathrm{H}_2$ gas at $300\,$K, $1\,$atm for 30 min.
(ii)~Press into a pellet ($\sim 5\,$mm diameter) at
$500\,$MPa.  (iii)~Mount in a cryostat with 4-point resistance
probe.  (iv)~Cool to $4\,$K, measure $R(T)$ on warming.
Expected: resistivity drop at $T \approx 31\,$K, zero resistance
below $\sim 28\,$K.  (v)~Confirm with SQUID magnetometry:
diamagnetic signal below $\Tc$.

The complete experimental protocol is given in Supplemental
Material~\cite{supplement}, Sec.~S7A.

\subsection{Prediction 2: Pressure-quench LaH$_{10}$ to
ambient conditions}

LaH$_{10}$ at $170\,$GPa has $\Tc = 260\,$K $= -13\,^\circ$C
with $\Rtau = 4.24$---far above unity.  The enormous $\Rtau$
surplus means that even if $\Icoop$ is substantially reduced
upon decompression, superconductivity may survive.

The key question is whether the $Fm\bar{3}m$ clathrate structure
of LaH$_{10}$ can be metastably retained at ambient pressure.
Deng and Chu~\cite{DengChu2026} demonstrated that
pressure-quenched HgBaCaCuO retains $\Tc = 151\,$K at
ambient pressure, proving the viability of the approach for
complex oxides.

\textbf{Proposed protocol.}---(i)~Synthesize LaH$_{10}$ at
$170\,$GPa, $2000\,$K in a diamond anvil cell
(DAC)~\cite{Drozdov2019}.  (ii)~Cool to $4\,$K under full
pressure.  (iii)~Rapidly decompress ($dP/dt > 10\,$GPa/s)
while maintaining $T < 20\,$K.  (iv)~Warm slowly, measuring
$R(T)$.

The rapid low-temperature quench exploits kinetic trapping:
at $4\,$K, hydrogen diffusion is negligible
($D_H \sim 10^{-30}\,$cm$^2$/s), so the clathrate cage
cannot decompose on experimental timescales.  Even partial
retention of the structure ($\Rtau$ reduced from 4.24 to
$\sim 1.5$) would yield $\Tc \sim 180$--$200\,$K at
ambient pressure---near $0\,^\circ$C.

Risk: LaH$_{10}$ may decompose via hydrogen escape through
grain boundaries even at low $T$.  Mitigation: encapsulate
in an amorphous boron nitride shell before decompression.

See Supplemental Material~\cite{supplement}, Sec.~S7B for the
complete DAC protocol.

\subsection{Prediction 3: Multi-channel enhancement in
La$_3$Ni$_2$O$_7$ $+$ H}

La$_3$Ni$_2$O$_7$ already achieves
$\Rtau^{\mathrm{multi}} = 1.47$ at $14\,$GPa through
3 pairing channels (phonon $+$ spin $+$ Ni $3d$ orbital).
Adding hydrogen via topotactic hydrogenation would introduce
a fourth channel: hydrogen vibrational modes coupling to the
Ni $3d$ electrons.

The expected $\Tc$ enhancement is
\begin{equation}\label{eq:Tc_enhanced}
  \frac{\Tc^{\mathrm{H}}}{\Tc^{(0)}} \approx
  \frac{\Rtau^{\mathrm{multi}}(4\,\text{ch})}
  {\Rtau^{\mathrm{multi}}(3\,\text{ch})}
  = \frac{4}{3}\,,
\end{equation}
predicting $\Tc^{\mathrm{H}} \approx 107 \pm 25\,$K at 14\,GPa,
or $\Tc \approx 84 \pm 20\,$K for the ambient-pressure film.

\subsection{Prediction 4: Eu$_2$H$_6$Ir
(4-channel candidate)}

Among the 897 screened compounds, Eu$_2$H$_6$Ir has the
highest $\Rtau^{\mathrm{multi}} = 0.38$ at ambient pressure,
with four identified channels:
(i)~phonon (H vibrational modes, $\omega_H \approx 100\,$meV);
(ii)~spin (Eu $4f^7$, $S = 7/2$, large local moment);
(iii)~topological (Ir $5d$ spin-orbit coupling,
potentially nontrivial $\mathbb{Z}_2$ invariant);
(iv)~orbital (Ir $5d$--Eu $5d$ hybridization).

Despite $\Rtau^{\mathrm{multi}} < 1$, this compound may
exhibit superconductivity at $\Tc \approx 10$--$30\,$K
via the multi-channel mechanism.  It has not been synthesized.

\textbf{Proposed synthesis.}---Arc-melt Eu $+$ Ir in a
1:0.5 stoichiometric ratio under argon.  Hydrogenate the
resulting Eu$_2$Ir ingot at $300\,$K, $50\,$bar H$_2$
for 24\,h.  Characterize structure by XRD; measure $R(T)$
down to $2\,$K.

%=======================================================================
\section{Connection to quantum processor experiments}
\label{sec:tuna9}
%=======================================================================

The $\tauasym$ framework has a direct experimental antecedent.
On the QuTech Tuna-9 superconducting processor~\cite{Huang2026exp},
a single Bell pair reduced $\tauasym$ from 0.47 to 0.05 (89\%
reduction).  This is the same physics as Cooper pairing:

\begin{center}
\begin{tabular}{lcc}
\toprule
& Bell pair (Tuna-9) & Cooper pair \\
\midrule
Entanglement & $\ln 2$ nats & $\ln 2$ nats (at $\bm{k}_F$) \\
Bare $\tauasym$ & 0.47 & $\taudec(T)$ \\
Assisted $\tauasym$ & 0.05 & 0 (below $\Tc$) \\
Recovery efficiency & 89\% & 100\% (macroscopic) \\
\bottomrule
\end{tabular}
\end{center}

The Cooper pair achieves 100\% recovery (vs.\ 89\% for Tuna-9)
because the BCS condensate involves $\sim 10^{23}$ entangled
pairs acting collectively, whereas Tuna-9 used a single pair
on a noisy channel.  Future quantum processor experiments with
tunable noise could directly simulate the $\Rtau$ criterion for
candidate materials~\cite{Huang2026exp}.

%=======================================================================
\section{Discussion}
\label{sec:discussion}
%=======================================================================

\subsection{What the $\Rtau$ criterion adds}

The $\Rtau$ screening criterion provides three capabilities
absent from existing approaches:

\emph{(i) Mechanism-independent comparison.}---Table~\ref{tab:known}
places conventional (Nb), two-gap (MgB$_2$), cuprate (YBCO),
hydride (H$_3$S), and nickelate (La$_3$Ni$_2$O$_7$)
superconductors on a single quantitative scale.
This is impossible with the McMillan formula alone, which
requires knowledge of the pairing mechanism.

\emph{(ii) High-throughput screening.}---The proxy
$\Rtau \approx 1.22\, \Tc\, \ThetaD / T^2$ requires only two
numbers per material and can be evaluated for entire databases.
Our screening of 897 compounds took $< 1$ CPU-hour.

\emph{(iii) Design guidance.}---The multi-channel
criterion~\eqref{eq:RT_multi} quantifies the advantage of
combining pairing mechanisms, providing a concrete target
for materials design.

\subsection{Relation to prior work on entanglement in superconductors}

Previous work has established that the BCS ground state contains
macroscopic entanglement~\cite{Vedral2004,Dunning2005}, that
Cooper pairs can be viewed as entangled
pairs~\cite{OhKim2004,Botero2004}, and that information-theoretic
measures detect the superconducting
transition~\cite{Ding2021}.  Amico et al.~\cite{Amico2008}
provide a comprehensive review of entanglement in
many-body systems, and Brand\~ao~\cite{Brandao2005} analyzed
entanglement quantification in mixed states relevant to
finite-temperature superconductors.  Most recently,
Huo et al.~\cite{Huo2024} applied the Petz recovery map to
diagnose quantum phases in topological systems.

Our contribution is distinct: we use the Petz recovery framework
not as a diagnostic but as a \emph{mechanism}---proposing that
the entanglement of Cooper pairs directly cancels the temporal
asymmetry of electron transport, yielding the $\Rtau$ screening
criterion.  This mechanistic role of entanglement is what
produces the quantitative material predictions absent in prior
information-theoretic analyses of superconductivity.

\subsection{Why room temperature is hard}

The screening result---no ambient-pressure single-channel
material among 897 hydrogen-metallic compounds reaches
$\Rtau = 1$---has a simple physical explanation.
At $T = 300\,$K, the phonon decoherence
$\taudec \approx \pi\lambda \cdot 300/\ThetaD$ grows linearly
with $T$.  To compensate, $\Icoop \propto \Delta(0) \propto
\ThetaD\, e^{-1/\lambda}$ must be exponentially large in
$1/\lambda$.  The product $\Tc\, \ThetaD$ is bounded from
above~\cite{Gao2025}: for Eliashberg phonon-mediated
superconductors at ambient pressure, $\Tc\, \ThetaD \lesssim
30{,}000\;\mathrm{K}^2$, which falls short of the single-channel
threshold $73{,}800\;\mathrm{K}^2$ by a factor of $\sim 2.5$.

The high-pressure hydrides circumvent this by achieving
$\ThetaD > 1000\,$K through the light hydrogen mass and
stiff covalent bonds.  The challenge of room-temperature
superconductivity at ambient pressure thus reduces to:
\emph{how to retain $\ThetaD > 1000\,$K without megabar
pressures?}  The $\Rtau$ framework suggests three routes:
metastable quenching, multi-channel pairing, or a new
mechanism.

\subsection{Limitations}

(i)~The proxy~\eqref{eq:Rtau_simple} assumes weak coupling
and neglects retardation effects.  For strong-coupling
materials ($\lambda > 2$), full Eliashberg calculations are
needed to validate $\Rtau$.

(ii)~The multi-channel enhancement assumes independent
(incoherent) addition.  Inter-channel coherence may enhance or
suppress $\Rtau^{\mathrm{multi}}$; the conditions for constructive
interference require symmetry analysis on a case-by-case basis.

(iii)~The Materials Project screening relies on estimated $\lambda$
values with $\sim 30$--$50\%$ uncertainty.  First-principles
electron-phonon calculations (e.g., via EPW~\cite{Ponce2016})
would sharpen the predictions but are computationally expensive
for all 897 compounds.

(iv)~The topological contribution $I_{\mathrm{topo}} = |\nu|\ln 2$
is temperature-independent by construction, but its magnitude
for specific materials requires knowledge of the band topology,
which is not available in the Materials Project for most entries.

%=======================================================================
\section{Conclusion}
\label{sec:conclusion}
%=======================================================================

We have developed the $\Rtau$ screening criterion for
superconductivity from the quantum information framework
$\tauasym = 1 - F$.  The criterion is mechanism-independent,
computationally inexpensive, and recovers the BCS gap equation
as a special case.  Screening 897 hydrogen-metallic compounds
reveals that ambient-pressure room-temperature superconductivity
via single-channel phonon pairing is extremely unlikely, but
multi-channel materials and pressure-quenched hydrides remain
viable.

While pressure quenching of known hydrides (LaH$_{10}$, H$_3$S)
provides the most direct route to near-$0\,^\circ$C
superconductivity, the long-term goal of ambient-pressure
room-temperature superconductivity requires discovering new
materials.  The $\Rtau$ framework provides three concrete
strategies: (1)~light-element scaffolds for maximum $\ThetaD$,
(2)~multi-channel $s$-wave pairing, and (3)~metastable phase
engineering.  We have identified three specific material families
(clathrate hydrides, nickelate hydrides, diamond-scaffold
hydrides) as priority targets for experimental synthesis and
first-principles screening.

Our most testable suggestion is that LaNi$_5$H$_6$---a
ubiquitous hydrogen storage alloy---may superconduct at
$\Tc \approx 31 \pm 12\,$K based on heuristic coupling
estimates.  This can be verified with standard cryogenic
equipment in any condensed matter laboratory.
If confirmed, it would validate the $\Rtau$ framework and
open a systematic search for higher-$\Tc$ multi-channel materials.

The deepest lesson is that a Cooper pair is a quantum error
correction code: the BCS condensate is a macroscopic ensemble
of entanglement-assisted recovery channels, each protecting
electronic quantum information against thermal
decoherence~\cite{Huang2026,Huang2026exp}.  Room-temperature
superconductivity requires codes that can correct errors at
a rate set by $300\,$K phonons---a formidable but perhaps not
impossible challenge.

\begin{acknowledgments}
The author thanks QuTech and the Quantum Inspire team for cloud
access to the Tuna-9 processor, and the Materials Project team
for the open database that enabled the screening study.
\end{acknowledgments}

%=======================================================================
\begin{thebibliography}{60}

\bibitem{supplement}
See Supplemental Material at [URL] for detailed derivations,
material-by-material calculations, experimental protocols,
and the $\Rtau$ calculator code.

\bibitem{Onnes1911}
H.~K.~Onnes,
The resistance of pure mercury at helium temperatures,
Commun.\ Phys.\ Lab.\ Univ.\ Leiden \textbf{12}, 120 (1911).

\bibitem{BCS1957}
J.~Bardeen, L.~N.~Cooper, and J.~R.~Schrieffer,
Theory of superconductivity,
Phys.\ Rev.\ \textbf{108}, 1175 (1957).

\bibitem{Bednorz1986}
J.~G.~Bednorz and K.~A.~M\"uller,
Possible high $T_c$ superconductivity in the Ba--La--Cu--O system,
Z.\ Phys.\ B \textbf{64}, 189 (1986).

\bibitem{Schilling1993}
A.~Schilling, M.~Cantoni, J.~D.~Guo, and H.~R.~Ott,
Superconductivity above 130\,K in the Hg--Ba--Ca--Cu--O system,
Nature \textbf{363}, 56 (1993).

\bibitem{Drozdov2015}
A.~P.~Drozdov, M.~I.~Eremets, I.~A.~Troyan, V.~Ksenofontov, and S.~I.~Shylin,
Conventional superconductivity at 203\,kelvin at high pressures in the sulfur hydride system,
Nature \textbf{525}, 73 (2015).

\bibitem{Somayazulu2019}
M.~Somayazulu \textit{et al.},
Evidence for superconductivity above 260\,K in lanthanum superhydride at megabar pressures,
Phys.\ Rev.\ Lett.\ \textbf{122}, 027001 (2019).

\bibitem{Sun2023}
H.~Sun \textit{et al.},
Signatures of superconductivity near 80\,K in a nickelate under high pressure,
Nature \textbf{621}, 493 (2023).

\bibitem{Ko2024}
E.~Ko \textit{et al.},
Signatures of ambient pressure superconductivity in thin film La$_3$Ni$_2$O$_7$,
Nature \textbf{638}, 935 (2025).

\bibitem{DengChu2026}
D.~Deng and C.~W.~Chu,
Pressure-quenched HgBaCaCuO with $T_c = 151\,$K at ambient pressure,
Preprint (2026).

\bibitem{McMillan1968}
W.~L.~McMillan,
Transition temperature of strong-coupled superconductors,
Phys.\ Rev.\ \textbf{167}, 331 (1968).

\bibitem{Allen1975}
P.~B.~Allen and R.~C.~Dynes,
Transition temperature of strong-coupled superconductors reanalyzed,
Phys.\ Rev.\ B \textbf{12}, 905 (1975).

\bibitem{Huang2026}
S.-K.~Huang,
Temporal asymmetry as Petz recovery failure,
Zenodo DOI: 10.5281/zenodo.18897853 (2026);
see also \url{https://github.com/akaiHuang/petz-recovery-unification}.

\bibitem{Huang2026exp}
S.-K.~Huang,
Measuring temporal asymmetry on a superconducting quantum processor:
From quantum erasure to the birth of time's arrow,
Preprint (2026).

\bibitem{Jain2013}
A.~Jain \textit{et al.},
Commentary: The Materials Project: A materials genome approach to accelerating materials innovation,
APL Mater.\ \textbf{1}, 011002 (2013).

\bibitem{Chen2019}
C.~Chen \textit{et al.},
A critical review of machine learning of energy materials,
Adv.\ Energy Mater.\ \textbf{10}, 1903242 (2020).

\bibitem{Stanev2018}
V.~Stanev \textit{et al.},
Machine learning modeling of superconducting critical temperature,
npj Comput.\ Mater.\ \textbf{4}, 29 (2018).

\bibitem{Petz1988}
D.~Petz,
Sufficiency of channels over von Neumann algebras,
Q.\ J.\ Math.\ \textbf{39}, 97 (1988).

\bibitem{Uhlmann1976}
A.~Uhlmann,
The ``transition probability'' in the state space of a *-algebra,
Rep.\ Math.\ Phys.\ \textbf{9}, 273 (1976).

\bibitem{Junge2018}
M.~Junge, R.~Renner, D.~Sutter, M.~M.~Wilde, and A.~Winter,
Universal recovery maps and approximate sufficiency of quantum relative entropy,
Ann.\ Henri Poincar\'e \textbf{19}, 2955 (2018).

\bibitem{Fawzi2015}
O.~Fawzi and R.~Renner,
Quantum conditional mutual information and approximate Markov chains,
Commun.\ Math.\ Phys.\ \textbf{340}, 575 (2015).

\bibitem{Devetak2005}
I.~Devetak and A.~Winter,
Distillation of secret key and entanglement from quantum states,
Proc.\ R.\ Soc.\ A \textbf{461}, 207 (2005).

\bibitem{Eliashberg1960}
G.~M.~Eliashberg,
Interactions between electrons and lattice vibrations in a superconductor,
Sov.\ Phys.\ JETP \textbf{11}, 696 (1960).

\bibitem{Choi2002}
H.~J.~Choi, D.~Roundy, H.~Sun, M.~L.~Cohen, and S.~G.~Louie,
The origin of the anomalous superconducting properties of MgB$_2$,
Nature \textbf{418}, 758 (2002).

\bibitem{Liu2017}
H.~Liu, I.~I.~Naumov, R.~Hoffmann, N.~W.~Ashcroft, and R.~J.~Hemley,
Potential high-$T_c$ superconducting lanthanum and yttrium hydrides at high pressure,
Proc.\ Natl.\ Acad.\ Sci.\ USA \textbf{114}, 6990 (2017).

\bibitem{Errea2020}
I.~Errea \textit{et al.},
Quantum crystal structure in the 250-kelvin superconducting lanthanum hydride,
Nature \textbf{578}, 66 (2020).

\bibitem{Gao2025}
M.~Gao \textit{et al.},
Upper bound on the critical temperature of superconductors from
first-principles phonon calculations,
Phys.\ Rev.\ B \textbf{111}, 024506 (2025).

\bibitem{Drozdov2019}
A.~P.~Drozdov \textit{et al.},
Superconductivity at 250\,K in lanthanum hydride under high pressures,
Nature \textbf{569}, 528 (2019).

\bibitem{VanVucht1970}
J.~H.~N.~van~Vucht, F.~A.~Kuijpers, and H.~C.~A.~M.~Bruning,
Reversible room-temperature absorption of large quantities of hydrogen by
intermetallic compounds,
Philips Res.\ Rep.\ \textbf{25}, 133 (1970).

\bibitem{Percheron-Guegan1985}
A.~Percheron-Gu\'egan, C.~Lartigue, J.~C.~Achard, P.~Germi, and F.~Tasset,
Neutron and X-ray diffraction profile analyses and structure of LaNi$_5$,
LaNi$_{5-x}$Al$_x$ and LaNi$_{5-x}$Mn$_x$ intermetallics and their hydrides,
J.\ Less-Common Met.\ \textbf{74}, 1 (1980).

\bibitem{Gupta1987}
M.~Gupta,
Electronic structure and electron-phonon coupling in LaNi$_5$,
J.\ Less-Common Met.\ \textbf{130}, 219 (1987).

\bibitem{Richter1981}
D.~Richter, R.~Hempelmann, and R.~C.~Bowman,
Dynamics of hydrogen in intermetallic hydrides,
in \textit{Hydrogen in Intermetallic Compounds II},
edited by L.~Schlapbach (Springer, Berlin, 1992), p.~97.

\bibitem{Paul-Boncour2004}
V.~Paul-Boncour \textit{et al.},
Structural and magnetic properties of ErFe$_2$H$_x$ and YFe$_2$H$_x$ compounds
($0 \leq x \leq 5$),
J.\ Solid State Chem.\ \textbf{142}, 120 (1999).

\bibitem{Ponce2016}
S.~Ponc\'e, E.~R.~Margine, C.~Verdi, and F.~Giustino,
EPW: Electron-phonon coupling, transport and superconducting properties
using maximally localized Wannier functions,
Comput.\ Phys.\ Commun.\ \textbf{209}, 116 (2016).

\bibitem{Anderson1987}
P.~W.~Anderson,
The resonating valence bond state in La$_2$CuO$_4$ and superconductivity,
Science \textbf{235}, 1196 (1987).

\bibitem{Scalapino1986}
D.~J.~Scalapino, E.~Loh, and J.~E.~Hirsch,
$d$-wave pairing near a spin-density-wave instability,
Phys.\ Rev.\ B \textbf{34}, 8190 (1986).

\bibitem{Moriya2000}
T.~Moriya and K.~Ueda,
Spin fluctuations and high temperature superconductivity,
Adv.\ Phys.\ \textbf{49}, 555 (2000).

\bibitem{Kontani2010}
H.~Kontani and S.~Onari,
Orbital-fluctuation-mediated superconductivity in iron pnictides:
Analysis of the five-orbital Hubbard--Holstein model,
Phys.\ Rev.\ Lett.\ \textbf{104}, 157001 (2010).

\bibitem{Fu2008}
L.~Fu and C.~L.~Kane,
Superconducting proximity effect and Majorana fermions at the surface
of a topological insulator,
Phys.\ Rev.\ Lett.\ \textbf{100}, 096407 (2008).

\bibitem{Sato2017}
M.~Sato and Y.~Ando,
Topological superconductors: A review,
Rep.\ Prog.\ Phys.\ \textbf{80}, 076501 (2017).

\bibitem{Ziman1960}
J.~M.~Ziman,
\textit{Electrons and Phonons} (Oxford University Press, Oxford, 1960).

\bibitem{ReebWolf2014}
D.~Reeb and M.~M.~Wolf,
An improved Landauer principle with finite-size corrections,
New J.\ Phys.\ \textbf{16}, 103011 (2014).

\bibitem{Tinkham2004}
M.~Tinkham,
\textit{Introduction to Superconductivity}, 2nd ed.\ (Dover, New York, 2004).

\bibitem{Vedral2004}
V.~Vedral,
High-temperature macroscopic entanglement,
New J.\ Phys.\ \textbf{6}, 102 (2004).

\bibitem{Zurek2019}
E.~Zurek and T.~Bi,
High-temperature superconductivity in alkaline and rare earth polyhydrides at high pressure:
A theoretical perspective,
J.\ Chem.\ Phys.\ \textbf{150}, 050901 (2019).

\bibitem{Mattis1958}
D.~C.~Mattis and J.~Bardeen,
Theory of the anomalous skin effect in normal and superconducting metals,
Phys.\ Rev.\ \textbf{111}, 412 (1958).

\bibitem{Boeri2022}
L.~Boeri \textit{et al.},
The 2021 room-temperature superconductivity roadmap,
J.\ Phys.\ Condens.\ Matter \textbf{34}, 183002 (2022).

\bibitem{Flores-Livas2020}
J.~A.~Flores-Livas \textit{et al.},
A perspective on conventional high-temperature superconductors at high pressure:
Methods and materials,
Phys.\ Rep.\ \textbf{856}, 1 (2020).

\bibitem{Pickard2020}
C.~J.~Pickard, I.~Errea, and M.~I.~Eremets,
Superconducting hydrides under pressure,
Annu.\ Rev.\ Condens.\ Matter Phys.\ \textbf{11}, 57 (2020).

\bibitem{Amico2008}
L.~Amico, R.~Fazio, A.~Osterloh, and V.~Vedral,
Entanglement in many-body systems,
Rev.\ Mod.\ Phys.\ \textbf{80}, 517 (2008).

\bibitem{Ding2021}
W.~Ding \textit{et al.},
Diagnosing quantum phase transitions via holographic entanglement entropy,
Proc.\ Natl.\ Acad.\ Sci.\ USA \textbf{118}, e2104114118 (2021).

\bibitem{OhKim2004}
S.~Oh and J.~Kim,
Entanglement of electron spins of noninteracting electron gases,
arXiv:quant-ph/0406013 (2004).

\bibitem{Dunning2005}
C.~Dunning, J.~Links, and H.-Q.~Zhou,
Ground-state entanglement of the BCS model,
Phys.\ Rev.\ Lett.\ \textbf{94}, 227002 (2005).

\bibitem{Botero2004}
A.~Botero,
Quantum information and the BCS ground state,
arXiv:quant-ph/0404176 (2004).

\bibitem{Brandao2005}
F.~G.~S.~L.~Brand\~ao,
Quantifying entanglement with witness operators,
Phys.\ Rev.\ A \textbf{72}, 022310 (2005).

\bibitem{Huo2024}
M.~Huo \textit{et al.},
Petz recovery map diagnostics of quantum phases,
Phys.\ Rev.\ B \textbf{110}, 195107 (2024).

\end{thebibliography}

\end{document}
